\documentclass[12pt,final]{article}
\usepackage{econ_paper_2}

\begin{document}
\setpaperheader{Model 2: Baseline New Keynesian Model with Capital}

\section{Model 2: Baseline with Capital}
\label{sec:model2-clean}

This document collects the equations for Model 2 in a compact, Dynare-oriented order. It retains the equations and timing in the original model and includes the equations inherited from Model 1.

\subsection{Variables}

The endogenous variables are itemized in the following order:
\begin{enumerate}[label=\arabic*.,leftmargin=2.4em]
  \item $r_t^k$: real rental rate of capital,
  \item $w_t$: real wage rate,
  \item $y_t$: aggregate output before price-adjustment costs,
  \item $n_t$: labor hours,
  \item $c_t$: household consumption,
  \item $i_t$: investment,
  \item $k_t$: capital available at the beginning of period $t$,
  \item $\pi_t=p_t/p_{t-1}$: gross inflation rate,
  \item $r_t$: gross nominal interest rate,
  \item $\Psi_t$: real marginal cost,
  \item $z_t$: level of technology.
\end{enumerate}

Additional variables and objects appearing in the model are
\begin{itemize}[leftmargin=2em]
  \item $b_t$: real value of a one-period nominal bond,
  \item $d_t$: real profits distributed by intermediate-goods firms,
  \item $y_t^{adj}$: output net of Rotemberg price-adjustment costs,
  \item $p_t(f)$ and $y_t(f)$: the price and output of intermediate firm $f$,
  \item $E_t[\cdot]$: expectation conditional on information available in period $t$,
  \item $\lambda_{t,k}$: the household's real stochastic discount factor from period $t$ to period $k$.
\end{itemize}

The exogenous innovation is
\begin{equation*}
  \upsilon_t \sim \mathbb{N}(0,\sigma_{\upsilon}^{2}),
\end{equation*}
which drives the technology process.

\subsection{Balanced-Growth Variable Classification}
\label{sec:bgp-variable-classification}

This subsection records the notation and trend classification that will govern
the balanced-growth rewrite.  It does not yet alter any equilibrium equation
below.  Until those equations are rewritten one block at a time, their lowercase
symbols retain the level meanings stated in the preceding variable list.

Following Eusepi and Preston (2011), let $X_t$ denote the labor-augmenting
stochastic trend.  Uppercase letters will denote nonstationary levels and
lowercase letters their stationary normalizations.  The intended classification
is:

\begin{table}[H]
  \centering
  \small
  \begin{tabular}{llll}
    \toprule
    Economic object & Level notation & Balanced-growth behavior
      & Planned stationary object \\
    \midrule
    Output & $Y_t$ & grows with $X_t$ & $y_t=Y_t/X_t$ \\
    Consumption & $C_t$ & grows with $X_t$ & $c_t=C_t/X_t$ \\
    Investment & $I_t$ & grows with $X_t$ & $i_t=I_t/X_t$ \\
    Real wage & $W_t$ & grows with $X_t$ & $w_t=W_t/X_t$ \\
    Capital stock & $K_t$ & predetermined at the start of $t$
      & $k_t=K_t/X_{t-1}$ \\
    Adjusted output & $Y_t^{adj}$ & grows with $X_t$
      & $y_t^{adj}=Y_t^{adj}/X_t$ \\
    Price-adjustment cost & $Adj_t$ & grows with $X_t$
      & $adj_t=Adj_t/X_t$ \\
    Real profits & $D_t$ & grows with $X_t$ & $d_t=D_t/X_t$ \\
    Real bond position & $B_t$ & grows with $X_t$ if nonzero
      & $b_t=B_t/X_t$ \\
    Hours & $N_t$ & stationary & $n_t=N_t$ \\
    Real capital rental rate & $R_t^k$ & stationary & $r_t^k=R_t^k$ \\
    Gross inflation & $\Pi_t$ & stationary & $\pi_t=\Pi_t$ \\
    Gross nominal rate & $R_t$ & stationary & $r_t=R_t$ \\
    Real marginal cost & $\Psi_t$ & stationary & $\psi_t=\Psi_t$ \\
    Relative firm price & $P_t(f)/P_t$ & stationary
      & $p_t(f)/p_t=P_t(f)/P_t$ \\
    Technology level & $X_t$ & stochastic trend & not stationary \\
    Gross technology growth & $\Gamma_t=X_t/X_{t-1}$ & stationary
      & $\gamma_t=\log\Gamma_t$ \\
    \bottomrule
  \end{tabular}
  \caption{Planned balanced-growth classification.  ``Stationary'' here means
  stationary around a constant balanced-growth steady state, not necessarily
  constant period by period.}
  \label{tab:bgp-variable-classification}
\end{table}

As in Eusepi and Preston, $K_t$ is predetermined capital available at the
beginning of period $t$, and $k_t=K_t/X_{t-1}$.  Investment chosen in period $t$
determines $K_{t+1}$.  This is a reindexing of the original document's
end-of-period convention, not a change in the underlying accumulation technology.

\subsection{Parameters}

\begin{description}[style=multiline,leftmargin=3.2cm,font=\normalfont]
  \item[$\beta$] household discount factor;
  \item[$\chi$] scale parameter on the disutility of labor;
  \item[$\eta$] inverse of the Frisch elasticity of labor supply, so the Frisch elasticity is $1/\eta$;
  \item[$\delta$] capital depreciation rate;
  \item[$\alpha$] cost share of capital in production;
  \item[$\theta$] elasticity of substitution between differentiated intermediate goods;
  \item[$\varphi$] Rotemberg price-adjustment cost coefficient;
  \item[$\bar{\pi}$] steady-state gross inflation rate;
  \item[$\phi_{\pi}$] monetary-policy response to inflation;
  \item[$\phi_y$] monetary-policy response to output;
  \item[$r^*$] intercept in the monetary-policy rule;
  \item[$\pi^*$] inflation target, with $\pi^*=\bar{\pi}$;
  \item[$y_t^*$] output target, equal either to steady-state output $\bar y$ or potential output $y_t^n$;
  \item[$\bar z$] steady-state technology;
  \item[$\rho_z$] persistence of technology, with $0\leq\rho_z<1$;
  \item[$\sigma_{\upsilon}$] standard deviation of the technology innovation;
  \item[$\mu$] markup term appearing in the paper's expression for potential output.
\end{description}

The paper's baseline calibration for the parameters used directly in Model 2 is
\begin{align*}
  \beta  & = 0.995,
         & \frac{1}{\eta}  & = 3,
         & \delta          & = 0.025,
         & \alpha          & = 0.33,   \\
  \theta & = 6,
         & \varphi         & = 59.11,
         & \bar\pi         & = 1.006,
         & \phi_\pi        & = 1.5,    \\
  \phi_y & = 0.1,
         & \bar z          & = 1,
         & \rho_z          & = 0.9,
         & \sigma_\upsilon & = 0.0025.
\end{align*}
The parameter $\chi$ is calibrated so that steady-state labor equals approximately $1/3$.

\section{Economic Primitives}

\paragraph{Household objective.}
The representative household chooses $\{c_t,i_t,n_t,b_t\}_{t=0}^{\infty}$ to maximize the preferences inherited from Model 1:
\begin{equation}
  E_0\sum_{t=0}^{\infty}\beta^t
  \left[
  \log c_t-\frac{\chi n_t^{1+\eta}}{1+\eta}
  \right].
  \label{eq:utility}
\end{equation}

\paragraph{Household budget constraint.}
Consumption, investment, and bond purchases equal labor income, capital income, the real payoff from previously purchased bonds, and firm profits:
\begin{equation}
  c_t+i_t+b_t
  =w_tn_t+r_t^k k_t
  +\frac{r_{t-1}b_{t-1}}{\pi_t}+d_t.
  \label{eq:budget}
\end{equation}

\paragraph{Capital accumulation.}
Next period's capital equals undepreciated beginning-of-period capital plus
current investment:
\begin{equation}
  k_{t+1}=(1-\delta)k_t+i_t.
  \label{eq:capital-law}
\end{equation}

\paragraph{Intermediate-goods production.}
Each intermediate firm combines beginning-of-period capital and current labor
using a Cobb--Douglas technology:
\begin{equation}
  y_t(f)=z_t k_t(f)^{\alpha}n_t(f)^{1-\alpha}.
  \label{eq:firm-production}
\end{equation}
In a symmetric equilibrium, this becomes
\begin{equation}
  y_t=z_t k_t^{\alpha}n_t^{1-\alpha}.
  \label{eq:aggregate-production}
\end{equation}

\paragraph{Technology process.}
Technology follows the stationary stochastic process specified in the original model:
\begin{equation}
  z_t=\bar z\left(\frac{z_{t-1}}{\bar z}\right)^{\rho_z}
  \exp(\upsilon_t).
  \label{eq:technology}
\end{equation}

\paragraph{Final-good aggregation.}
The representative final-goods firm combines differentiated intermediate goods with a Dixit--Stiglitz aggregator:
\begin{equation}
  y_t
  \equiv
  \left[
    \int_0^1 y_t(f)^{(\theta-1)/\theta}\,df
    \right]^{\theta/(\theta-1)}.
  \label{eq:final-good}
\end{equation}

\paragraph{Demand for an intermediate good.}
Cost minimization by the final-goods firm yields demand for firm $f$'s output:
\begin{equation}
  y_t(f)=\left(\frac{p_t(f)}{p_t}\right)^{-\theta}y_t.
  \label{eq:firm-demand}
\end{equation}

\paragraph{Aggregate price index.}
The price of the final good is the corresponding Dixit--Stiglitz price index:
\begin{equation}
  p_t
  =\left[
    \int_0^1 p_t(f)^{1-\theta}\,df
    \right]^{1/(1-\theta)}.
  \label{eq:price-index}
\end{equation}

\paragraph{Rotemberg price-adjustment cost.}
Each intermediate firm loses output when it changes its price relative to steady-state inflation:
\begin{equation}
  adj_t(f)
  =\frac{\varphi}{2}
  \left[
    \frac{p_t(f)}{\bar\pi p_{t-1}(f)}-1
    \right]^2y_t.
  \label{eq:adjustment-cost-firm}
\end{equation}
In a symmetric equilibrium,
\begin{equation}
  adj_t
  =\frac{\varphi}{2}
  \left(\frac{\pi_t}{\bar\pi}-1\right)^2y_t.
  \label{eq:adjustment-cost}
\end{equation}

\paragraph{Intermediate-firm profits.}
The original Model 2 section does not restate the profit equation after adding capital. For a Dynare implementation, profits need not be introduced as a separate equilibrium variable because bond-market clearing and the aggregate resource constraint eliminate them from the compact system.

\paragraph{Pricing kernel.}
The household's one-period real stochastic discount factor is
\begin{equation}
  \lambda_{t,t+1}=\beta\frac{c_t}{c_{t+1}},
  \label{eq:pricing-kernel-one}
\end{equation}
with
\begin{equation}
  \lambda_{t,t}\equiv1,
  \qquad
  \lambda_{t,k}\equiv
  \prod_{j=t+1}^{k}\lambda_{j-1,j}.
  \label{eq:pricing-kernel-many}
\end{equation}

\section{Equilibrium Equation Block}

The following equations are the compact equilibrium system most directly corresponding to a Dynare model.

\paragraph{Household labor-supply condition.}
The real wage equals the household's marginal rate of substitution between consumption and labor:
\begin{equation}
  w_t=\chi n_t^{\eta}c_t.
  \label{eq:labor-foc}
\end{equation}

\paragraph{Nominal-bond Euler equation.}
The household is indifferent at the margin between current consumption and purchasing a one-period nominal bond:
\begin{equation}
  1
  =r_tE_t\left[
    \beta\frac{c_t}{c_{t+1}}\frac{1}{\pi_{t+1}}
    \right].
  \label{eq:bond-euler}
\end{equation}

\paragraph{Capital Euler equation.}
The marginal utility cost of saving through capital equals the expected discounted marginal utility of its rental return and undepreciated value:
\begin{equation}
  \frac{1}{c_t}
  =E_t\left[
    \beta\frac{1}{c_{t+1}}
    \left(r_{t+1}^k+1-\delta\right)
    \right].
  \label{eq:capital-euler}
\end{equation}

\paragraph{Capital-demand condition.}
The real rental rate of capital equals capital's marginal revenue product:
\begin{equation}
  r_t^k=\alpha\Psi_t\frac{y_t}{k_t}.
  \label{eq:capital-foc}
\end{equation}

\paragraph{Labor-demand condition.}
The real wage equals labor's marginal revenue product:
\begin{equation}
  w_t=(1-\alpha)\Psi_t\frac{y_t}{n_t}.
  \label{eq:firm-labor-foc}
\end{equation}

\paragraph{Aggregate production.}
Symmetric intermediate firms produce aggregate output using beginning-of-period
capital, current labor, and current technology:
\begin{equation}
  y_t=z_t k_t^{\alpha}n_t^{1-\alpha}.
  \label{eq:production-block}
\end{equation}

\paragraph{Rotemberg pricing condition.}
A price change today balances the current markup incentive against the expected value of future price adjustment:
\begin{equation}
  \varphi
  \left(\frac{\pi_t}{\bar\pi}-1\right)
  \frac{\pi_t}{\bar\pi}
  =(1-\theta)+\theta\Psi_t
  +\varphi E_t\left[
    \lambda_{t,t+1}
    \left(\frac{\pi_{t+1}}{\bar\pi}-1\right)
    \frac{\pi_{t+1}}{\bar\pi}
    \frac{y_{t+1}}{y_t}
    \right].
  \label{eq:pricing-foc}
\end{equation}

\paragraph{Monetary-policy rule.}
The central bank sets the gross nominal interest rate in response to inflation and output, subject to the zero lower bound:
\begin{equation}
  r_t
  =\max\left\{
  1,
  r^*\left(\frac{\pi_t}{\pi^*}\right)^{\phi_\pi}
  \left(\frac{y_t}{y_t^*}\right)^{\phi_y}
  \right\}.
  \label{eq:policy-rule}
\end{equation}
The inflation target is $\pi^*=\bar\pi$. The output target is either
\begin{equation}
  y_t^*=\bar y
  \label{eq:ss-output-target}
\end{equation}
or potential output,
% TODO: Verify this potential-output target for Model 2. The expression below is
% inherited from the labor-only Model 1 and does not depend on the predetermined
% capital stock. Either derive flexible-price output for the capital model or
% clarify that y_t^n is an externally specified policy benchmark rather than
% Model 2's actual output when \varphi=0.
\begin{equation}
  y_t^*=y_t^n=(\chi\mu)^{-1/(1+\eta)}z_t.
  \label{eq:potential-output-target}
\end{equation}
The paper also considers $\phi_y=0$.

\paragraph{Capital law of motion.}
Current investment determines next period's capital stock:
\begin{equation}
  k_{t+1}=(1-\delta)k_t+i_t.
  \label{eq:capital-block}
\end{equation}

\paragraph{Resource constraint.}
Consumption and investment exhaust output net of Rotemberg price-adjustment costs:
\begin{equation}
  c_t+i_t=y_t^{adj}.
  \label{eq:resource-adjusted}
\end{equation}
Adjusted output is defined by
\begin{equation}
  y_t^{adj}=y_t-adj_t
  =y_t-
  \frac{\varphi}{2}
  \left(\frac{\pi_t}{\bar\pi}-1\right)^2y_t.
  \label{eq:adjusted-output}
\end{equation}
Equivalently, without introducing $y_t^{adj}$ as a separate variable,
\begin{equation}
  c_t+i_t
  =y_t-
  \frac{\varphi}{2}
  \left(\frac{\pi_t}{\bar\pi}-1\right)^2y_t.
  \label{eq:resource-direct}
\end{equation}

\paragraph{Technology process.}
Technology evolves according to its stationary autoregressive process:
\begin{equation}
  z_t=\bar z\left(\frac{z_{t-1}}{\bar z}\right)^{\rho_z}
  \exp(\upsilon_t).
  \label{eq:technology-block}
\end{equation}

\paragraph{Bond-market clearing.}
Because the representative household's bond position is in zero net supply,
\begin{equation}
  b_t=0.
  \label{eq:bond-clearing}
\end{equation}

\section{Compact Dynare Ordering}

Using the eleven endogenous variables listed above, a compact equation order is:
\begin{enumerate}[label=\arabic*.,leftmargin=2.4em]
  \item capital-demand condition, \eqref{eq:capital-foc};
  \item household labor-supply condition, \eqref{eq:labor-foc};
  \item aggregate production, \eqref{eq:production-block};
  \item labor-demand condition, \eqref{eq:firm-labor-foc};
  \item capital Euler equation, \eqref{eq:capital-euler};
  \item direct resource constraint, \eqref{eq:resource-direct};
  \item capital law of motion, \eqref{eq:capital-block};
  \item Rotemberg pricing condition, \eqref{eq:pricing-foc};
  \item nominal-bond Euler equation, \eqref{eq:bond-euler};
  \item monetary-policy rule, \eqref{eq:policy-rule};
  \item technology process, \eqref{eq:technology-block}.
\end{enumerate}

This ordering is only organizational: it does not alter the model's equilibrium conditions.

\end{document}
